\chapter{4. Volume, injectivity radius and injectivity angle}
\source{cheeger-gromov-taylor:page-0028,cheeger-gromov-taylor:page-0029,cheeger-gromov-taylor:page-0030,cheeger-gromov-taylor:page-0031,cheeger-gromov-taylor:page-0032,cheeger-gromov-taylor:page-0033,cheeger-gromov-taylor:page-0034,cheeger-gromov-taylor:page-0035,cheeger-gromov-taylor:page-0036,cheeger-gromov-taylor:page-0037,cheeger-gromov-taylor:page-0038,cheeger-gromov-taylor:page-0039,cheeger-gromov-taylor:page-0040}

In this section we prove the geometric estimates used in the previous sections.

We begin by discussing the volume. Let $x\in M^n$, and $A(r,\theta)$,
$\mathcal{Q}^H(r)$, $V_r^H$ be as defined after (1.7) and in (1.16), (1.17).
Suppose $\Ric_M\geq (n-1)H$ and regard $A(r,\theta)$ as a function on the tangent
space $T_xM$. If $H>0$, we restrict attention to $B_{\pi/\sqrt{H}}(0)\subset T_xM$.
Let $U\subset T_xM$ denote the interior of the cut locus in the tangent space.
The proof of the standard volume comparison (see [3, pp.\ 253--257]) shows that
on $U$,
\begin{equation}
\tag{4.1}
\frac{A(r,\theta)}{\mathcal{Q}^H(r)}\downarrow,
\end{equation}
where $\downarrow$ indicates a decreasing function. If $S_r(0)\subset T_xM$ denotes
the sphere of radius $r$ about the origin, (4.1) implies
\begin{equation}
\tag{4.2}
\frac{\int_{S_r(0)\cap U} A(r,\theta)}{\int_{S_r(0)} \mathcal{Q}^H(r)}\downarrow;
\end{equation}
the possible existence of the cut points only helps matters. In general, if
$f(r)$, $g(r)$ are positive functions such that $\frac{f}{g}\downarrow$, then it follows that
\begin{equation}
\tag{4.3}
\frac{\int_0^r f(s)}{\int_0^r g(s)}\downarrow.
\end{equation}
To see this, set $f/g=h$. Then $h\downarrow$ and
\begin{equation}
\tag{4.4}
\int_0^r f \int_r^R g
= \int_0^r gh \int_r^R g
\geq \left[\int_0^r g\right] h(r)\int_r^R g
\geq \left[\int_0^r g\right]\int_r^R gh
= \int_0^r g \int_r^R f.
\end{equation}
Thus, equivalently,
\begin{equation}
\tag{4.5}
\int_0^r f \int_0^R g
= \int_0^r f \int_0^r g + \int_0^r f \int_r^R g
\geq \int_0^r f \int_0^r g + \int_0^r g \int_r^R f
= \int_0^r g \int_0^R f,
\end{equation}
which gives (4.3). From (4.2) and (4.3) we obtain the following relative volume
estimates which were emphasized in [21], [23].

\begin{proposition}[Relative volume estimates]
\label{prop:relative-volume-estimates}
\dcref{Proposition 4.1}
\group{section-4}
\source{cheeger-gromov-taylor:page-0028,cheeger-gromov-taylor:page-0029}
Let $M^n$ be complete and $\Ric_M\geq (n-1)H$. Let $p,x\in M^n$ and
$\rho(p,x)=\rho$. Then
\begin{equation}
\tag{i}
\frac{V_r(p)}{V_r^H}\downarrow,
\end{equation}
or equivalently, for $r_1 \le r_2 \le r_3$,
\begin{equation}
\tag{ii}
\frac{V_{r_3}(p)-V_{r_2}(p)}{V_{r_3}^H-V_{r_2}^H}\le \frac{V_{r_1}(p)}{V_{r_1}^H}.
\end{equation}
Moreover,
\begin{equation}
\tag{iii}
V_{r_2}(p)\frac{V_{r_1}^H}{V_{\rho+r_2}^H}\le V_{r_1}(x).
\end{equation}
If $r_2+r<\rho$, then
\begin{equation}
\tag{iv}
V_{r_2}(p)\frac{V_r^H}{V_{\rho+r_2}^H-V_{\rho-r_2}^H}\le V_r(x).
\end{equation}
\end{proposition}

\begin{proof}
\source{cheeger-gromov-taylor:page-0029}
(i) and (ii). Let
\begin{equation}
\tag{4.10}
f=\int_{S_r^{n-1}(0)\cap U}A(r,\theta), \qquad g=\int_{S_r^{n-1}(0)}\mathcal{Q}^H(r).
\end{equation}
Then (i) and (ii) follow from (4.2), (4.3).

(iii) We have
\begin{equation}
\tag{4.11}
V_{r_2}(p)\le V_{\rho+r_2}(x)\le V_{\rho+r_2}^H\frac{V_{r_1}(x)}{V_{r_1}^H}.
\end{equation}

(iv) Using (ii) we obtain
\begin{equation}
\tag{4.12}
V_{r_2}(p)\le V_{\rho+r_2}(x)-V_{\rho-r_2}(x)\le \left(V_{\rho+r_2}^H-V_{\rho-r_2}^H\right)\frac{V_r(x)}{V_r^H}.
\end{equation}
\end{proof}

Relative volume estimates are implicit in the proof of the more familiar volume
comparison $V_r(p)\le V_r^H$, and have, in fact, been used previously (compare
e.g.\ [5]). However, recent applications in [22], [23] indicate that their
significance was not fully appreciated.

Note that it follows immediately from Proposition 4.1(iv) that if
$\Ric_M\ge 0$, $d>2r$, then
\begin{equation}
\tag{4.13}
\frac{(d-2r)^n}{d^n-(d-2r)^n}V_r(p)\le V_d(p).
\end{equation}

Thus we recover the result of Calabi and Yau [35] that $V_d(p)$ grows linearly
as $d\to\infty$. At the end of this section we will give a sharp generalization
of (4.13) to manifolds whose Ricci curvature decays no slower than
$c/[\rho(p,x)^2]$ outside some $B_{r_0}(p)$.

\begin{remark}[Relative volume comparison for geodesically star-shaped sets]
\label{rem:relative-volume-star-shaped}
\dcref{Remark 4.1}
\group{section-4}
\source{cheeger-gromov-taylor:page-0029}
Let $S\subset M^n$ be any set with the property that if
$q\in S$ is not on the cut locus of $p$, then the unique minimal geodesic from $p$
to $q$ is contained in $S$. Clearly, Proposition 4.1(i) and (ii) continue to
hold with $V_r(p)$ replaced by $V(S\cap B_r(p))$, as do the corresponding
generalizations of (iii) and (iv).
\end{remark}

\begin{remark}[Ricci models]
\label{rem:ricci-models}
\dcref{Remark 4.2}
\group{section-4}
\source{cheeger-gromov-taylor:page-0030}
Let $M^n$ be a Ricci model in the sense of [9]. Thus $M^n$ is a metric ball
(possibly of infinite radius), with metric $dr^2+f^2(r)g$, where $f(r)\sim
r-Kr^3/6+\cdots$ for some $K$. If at distance $r$ from $p$, then we have
\[
\Ric_{M^n}\ge -(n-1)f''(r)/f(r),
\]
Remark 4.1 still applies with $V_r^H$ replaced by $f^{n-1}(r)$.
\end{remark}

Remarks 4.1 and 4.2 will be used without further comment below.
We now consider the refined injectivity angle $\wtomega(H,\epsilon,r,x)$
introduced in \S2 just before Proposition 2.3. We retain the assumptions of
Proposition 4.1.

\begin{proposition}[Lower estimates for the refined injectivity angle]
\label{prop:refined-injectivity-angle-lower-estimates}
\dcref{Proposition 4.2}
\group{section-4}
\source{cheeger-gromov-taylor:page-0030,cheeger-gromov-taylor:page-0031}
\textup{(i)} Given $\epsilon$, $r_2$, $r$, we have
\begin{equation}\tag{4.14}
\wtomega(H,\epsilon,r,x)\ge \frac{1}{1-\epsilon}
\left[\frac{V_{r_2}(p)-V_r^H}{V^H_{\rho+r_2}-V_r^H}-\epsilon\right],
\end{equation}
where the right-hand side is $>0$, if $\epsilon,r$ are sufficiently small.

\textup{(ii)} If moreover $r_2<\rho$, then
\begin{equation}\tag{4.15}
\wtomega(H,\epsilon,\rho-r_2,x)\ge \frac{1}{1-\epsilon}
\left[\frac{V_{r_2}(p)}{V^H_{\rho+r_2}-V^H_{\rho-r_2}}-\epsilon\right].
\end{equation}

\textup{(iii)} Thus if $r(H,V)$ is defined by $\frac{V}{2}=V^H_{r(H,V)}$, then
\begin{equation}\tag{4.16}
\wtomega\!\left(H,\frac{V_{r_2}(p)}{2(2V^H_{\rho+r_2}-V_{r_2}(p))},r(H,V_{r_2}(p)),x\right)
\ge \frac{V_{r_2}(p)}{4V^H_{\rho+r_2}-3V_{r_2}(p)}.
\end{equation}

\textup{(iv)} If moreover $r_2<\rho$, then
\[
\wtomega\!\left(H,\frac{V_{r_2}(p)}{4(V^H_{\rho+r_2}-V^H_{\rho-r_2})},\rho-r_2,x\right)
\ge \frac{V_{r_2}(p)}{4(V^H_{\rho+r_2}-V_{\rho-r_2}^H)-V_{r_2}(p)}.
\]
\end{proposition}

\begin{proof}
\source{cheeger-gromov-taylor:page-0030,cheeger-gromov-taylor:page-0031}
\textup{(i)} Since $V_r(x)\le V_r^H$ and $B_{r_2}(p)\subset B_{\rho+r_2}(x)$, we have
\begin{equation}\tag{4.17}
V_{r_2}(p)-V_r^H\le V_{\rho+r_2}(x)-V_r(x).
\end{equation}
Using (4.2) and integrating in polar coordinates, we easily get
\begin{equation}\tag{4.18}
V_{\rho+r_2}(x)-V_r(x)\le
\bigl(V^H_{\rho+r_2}-V_r^H\bigr)\bigl[\wtomega(H,\epsilon,r,x)+\epsilon(1-\wtomega(H,\epsilon,r,x))\bigr],
\end{equation}
from which (i) follows.

\textup{(ii)} This follows as in (i), by noting that
\begin{equation}
\tag{4.19}
V_{r_2}(p) \leq V_{\rho+r_2}(x) - V_{\rho-r_2}(x).
\end{equation}

\textup{(iii)} and (iv) follow immediately from (i) and (ii).
\end{proof}

We now discuss the injectivity radius. Let $M^n$ be a complete Riemannian
manifold and $p \in M^n$. Let $\gamma$ be a geodesic loop on $p$ of shortest
length $L[\gamma] = 2l$. Let $K_M \leq K$ on $B_l(p)$. Then by a result of
Klingenberg (see e.g.\ [7]), the injectivity radius of the exponential map at $p$
is bounded below by $\min(l, \pi/\sqrt{K})$; if $K \leq 0$, we interpret
$\pi/\sqrt{K}$ as $\infty$.

In [6] it was shown that there exists a constant $c_n(d, V, H) > 0$ with the
following properties. Let $M^n$ be \textit{compact}. Suppose the diameter
$d(M^n)$ and the volume $V(M^n)$ satisfy $d(M^n) \leq d$, $V(M^n) \geq V$, and
that $K_M \geq H$. Then every \textit{smooth} closed geodesic on $M^n$ has length
$\geq c_n(d, V, H)$. For compact mani- folds Klingenberg showed in addition that
the injectivity radius at all points is bounded below by the minimum of
$\pi/\sqrt{K}$ (where $K_M \leq K$) and one half the length of a \textit{smooth}
closed geodesic. Thus it follows that for all $p \in M^n$,
\begin{equation}
\tag{4.20}
i(p) \geq \min\left(\frac{1}{2} c_n(d, V, H), \pi/\sqrt{K}\right).
\end{equation}

The method of [6] depends on Toponogov's theorem. It leads to a sharp constant
in certain situations, and no part of the hypothesis can be removed unless
something further is added; see e.g.\ [7] for further results of various
authors in more special situations. However, in [24] Heintze and Karcher gave
a new derivation. Their method was to estimate the volume of a tube around a
smooth closed geodesic. This led to a better constant in most cases and did not
require the use of Toponogov's Theorem. Whether the lower bound
$K_M \geq H$ can be replaced by $\Ric_M \geq (n-1)H$ in the above estimates,
is still an open question.

As mentioned in the introduction, in [10] a relative bound on the behavior of
the injectivity radius $i(x)$ was derived, which did not require the global
hypothesis of compactness. In particular, the authors worked with geodesic
loops which were not necessarily smooth, but they assumed a lower bound on
$i(p)$ for some $p$. The estimate which we now give replaces this assumption by
a lower bound on the volume $V_r(p)$ for some $r$, $p$. Thus we get a purely
local generalization of the results of [6], [24], in so far as they pertain to
the injectivity radius. However, the method by its very nature depends on
knowing a lower bound for the distance to the conjugate locus. Without this
assumption, no information on the length of a closed geodesic $\gamma$ is
obtained, even if $\gamma$ is smooth.

To emphasize the local nature of our result, we now let $B_r(p)$ be a metric
ball in a Riemannian manifold such that for $r' < r$, $\overline{B_{r'}(p)}$ is
compact. Assume that on $B_r(p)$, $K_M \le K$ and $r \le \pi/\sqrt{K}$, ($r$
arbitrary if $K \le 0$). We let $V_r^0(p)$ denote the volume of $B_r(0) \subset
T_pM$ with respect to the pulled back metric.

\begin{theorem}[Local injectivity-radius bound]
\label{thm:local-injectivity-radius-bound}
\dcref{Theorem 4.3}
\group{section-4}
\source{cheeger-gromov-taylor:page-0032,cheeger-gromov-taylor:page-0033}
Let $B_r(p)$ be as above. If $\gamma$ is a geodesic loop on $p$ of length
$2l$, and $r_0+2s \le r$, $r_0 \le r/4$, then
\begin{equation}\tag{4.21}
l \ge \frac{r_0}{2}\,\frac{1}{1+V^0_{r_0+s}(p)/V_s(p)}.
\end{equation}
If in addition $H < K_M \le K$, then
\begin{equation}\tag{4.22}
l \ge \frac{r_0}{2}\,\frac{1}{1+V^H_{r_0+s}(p)/V_s(p)}.
\end{equation}
\end{theorem}

To prepare for the proof of Theorem 4.3, we begin with some elementary
observations.

\begin{enumerate}
\item[(i)] By the Gauss Lemma, every curve $c$ of length $L[c] < r$ with
$c(0)=p$ has a unique lift to a curve $\tilde c \subset B_r(0) \subset T_pM$ with
$\tilde c(0)=0$. In particular, a geodesic segment $\gamma$ of length $< r$ with
$\gamma(0)=p$ has a unique radial lift. Thus points of $B_r(0)$ can be
identified uniquely with such $\gamma$. With this identification, the projection
$\exp_p : B_r(0) \to B_r(p)$ is given by $\gamma \mapsto e(\gamma)$, where
$e(\gamma)$ denotes the endpoint of $\gamma$.

\item[(ii)] Let $c_1,c_2$ be piecewise smooth curves, and write $c_1 \sim_A c_2$
if $c_1$ and $c_2$ are homotopic over curves of length $\le A$ keeping endpoints
fixed. If $L[c]=l<r$, clearly $\tilde c \sim_l \tilde\gamma$ with $\tilde\gamma$
the unique radial geodesic such that $e(\tilde c)=e(\tilde\gamma)$. Hence
$c \sim_l \gamma$.

\item[(iii)] Since $\exp_p \vert_{B_r(0)}$ is nonsingular, if $\gamma_1,\gamma_2$
are distinct radial geodesics with $L[\gamma_i]<r$, then by a standard argument
on lifting homotopies, $\gamma_1,\gamma_2$ do not satisfy $\gamma_1 \sim_r
\gamma_2$. In particular, the geodesic segment $\gamma$ in (ii) is the unique
geodesic segment with $c \sim \gamma$. We write $[c]$ for $\gamma$, and
$[p]$ for the constant loop on $p$.
\end{enumerate}

Let $\theta_1,\theta_2$ be closed curves on $p$, and let $c$ be an arc from $p$
with $L[\theta_i]\le 2l$, $L[c]=s$ and $2(l+s)<r$. Let $\theta_i \cup c$ denote
the arc from $p$ which is equal to $\theta_i$ followed by $c$.

\begin{lemma}[Concatenation preserves radial classes]
\label{lem:radial-class-concatenation}
\dcref{Lemma 4.4}
\group{section-4}
\source{cheeger-gromov-taylor:page-0032}
If $2(l+s)<r$ and $[\theta_1\cup c]=[\theta_2\cup c]$, then $[\theta_1]=[\theta_2]$.
\end{lemma}

\begin{proof}
\source{cheeger-gromov-taylor:page-0032}
If $[\theta_1\cup c]=[\theta_2\cup c]$, then $\theta_1\cup c \sim_{2(l+s)}
\theta_2\cup c$. Thus
\[
\theta_1 \sim_{2(l+s)} \theta_1\cup c\cup -c \sim_{2(l+s)}
\theta_2\cup c\cup -c \sim_{2(l+s)} \theta_2.
\]
By (iii) above $[\theta_1]=[\theta_2]$.
\end{proof}

Let $\#(q,a)$ denote the number of distinct inverse images of $q \in B_r(p)$ in
$B_a(0)$, $a < r$.

\begin{lemma}[Even covering lemma]
\label{lem:even-covering}
\dcref{Lemma 4.5}
\group{section-4}
\source{cheeger-gromov-taylor:page-0033}
If $r_0 + 2s < r/4$ and $q \in B_s(p)$, then
\[
\#(q,3r_0) \ge \#(p,r_0).
\]
\end{lemma}

\begin{proof}
\source{cheeger-gromov-taylor:page-0033}
By (i) above, $\#(p,r_0)=m$, where $\gamma_1 \cdots \gamma_m$ are the distinct
geodesic loops on $p$ of length $< r_0$. Let $\sigma$ be a minimal geodesic from
$p$ to $q$. By Lemma 4.4, $[\gamma_i \cup \sigma]$, $i=1,\cdots,m$, represent
distinct inverse images of $q$ in $B_{r_0+s}(p)$.
\end{proof}

Let $\theta$ be a loop on $p$, and $\sigma$ a geodesic arc from $p$ with
$L(\theta)=2l$, $L(\sigma)=s$ and $2(l+s)<r$. The correspondence
$\sigma \mapsto [\theta \cup \sigma]$ defines a map $\pi_{[\theta]}: B_s(0) \to
B_{2l+s}(0)$ which clearly depends only on $[\theta]$. It is also obvious that
$\exp_p \pi_{[\theta]}(x)=\exp_p x$ for $x \in B_s(0)$. Hence $\pi_{[\theta]}$ is
an isometry from $B_s(0)$ to $\pi_{[\theta]}(B_s(0))$ with respect to the pulled
back metric on $B_{2l+s}(0)$. It follows immediately from Lemma 4.4 that
$\pi_{[\theta]}$ has no fixed points unless $[\theta]=p$.

Set $i\gamma=\gamma\cup \cdots \cup \gamma$ and $L(\gamma)=2l$.

\begin{lemma}[Iterated geodesic loops are distinct]
\label{lem:iterated-geodesic-loops}
\dcref{Lemma 4.6}
\group{section-4}
\source{cheeger-gromov-taylor:page-0033}
Let $\gamma$ be a geodesic loop on $p$ with $2l \cdot N < \frac14 r$. Then
$[\gamma], [2\gamma], \cdots, [N\gamma]$ are all distinct.
\end{lemma}

\begin{proof}
\source{cheeger-gromov-taylor:page-0033}
Since $(i+j)\gamma=i\gamma \cup j\gamma$, if
$[(i+j)\gamma]=[i\gamma \cup j\gamma]=[j\gamma]=[j\gamma \cup p]$, then
$[i\gamma]=p$ by Lemma 4.1. Regard $[p],[\gamma],\cdots,[(i-1)\gamma]$ as points
in the strictly convex ball $B_{2il}(0)$. Then $\pi_{[\gamma]}: B_{2il}(0) \to
B_{2(i+1)l}(0)$ preserving distance. Hence $\pi_{[\gamma]}$ preserves the unique
center of gravity of the set $\{[p],[\gamma],\cdots,[(i-1)\gamma]\}$. By
definition the center of gravity is the unique minimum of the function
$y \mapsto \sum_j \rho^2([j\gamma],y)$. It exists because $2il < \frac12
\pi/\sqrt{K}$ implies that $B_{2il}(0)$ is convex. This contradicts the fact that
$\pi_\gamma$ has no fixed points.
\end{proof}

\begin{proof}[Proof of Theorem 4.3.]
\source{cheeger-gromov-taylor:page-0033}
Consider a geodesic loop $\gamma$ with
$L(\gamma)=2l<r_0$. Let $N=\lfloor r_0/2l \rfloor$. By Lemma 4.6, $p$ has at
least $N$ inverse images in $B_{r_0}(0)$. By Lemma 4.5, every point in
$B_s(p)$ has at least $N$ inverse images in $B_{r_0+s}(0)$. Since $\exp_p$ is
nonsingular and orientation preserving, $N \cdot V_s(p) \le V^0_{r_0+s}(p)$.
Hence $1/\lfloor r_0/2l \rfloor \le V(B_s(p))/V(B_{r_0+s}(0))$, and (4.20)
follows. Then (4.21) follows from the standard volume comparison.
\end{proof}

If we combine Theorem 4.3 with Proposition 4.1, we immediately obtain the
following lower bound for the injectivity radius.

\begin{theorem}[Injectivity radius from relative volume]
\label{thm:injectivity-radius-relative-volume}
\dcref{Theorem 4.7}
\group{section-4}
\source{cheeger-gromov-taylor:page-0033,cheeger-gromov-taylor:page-0034}
Let $M^n$ be complete with $H \le K_M \le K$.
Let $\rho=\rho(p,x)$, and fix $r, r_0, s$, with $r_0 + 2s < \pi/\sqrt{k}$,
$r_0 \le \pi/4\sqrt{k}$.
\[
\text{(i)}
\]
\begin{equation}
\tag{4.23}
i(x) \ge \frac{r_0}{2}\,\frac{1}{1+\left(V^H_{r_0+s}/V_r(p)\right)\left(V^H_{\rho+r}/V^H_s\right)}.
\end{equation}
\[
\text{(ii)}
\]
Moreover, if $r+s<\rho$, then
\begin{equation}
\tag{4.24}
i(x)\ge \frac{r_0}{2}\,\frac{1}{1+\bigl(V^{H}_{r_0+s}/V_r(p)\bigr)\bigl(V^{H}_{\rho+r}-V^{H}_{\rho-r}\bigr)/V^{H}_s}.
\end{equation}
\end{theorem}

As an application of the relative volume estimates proved at the beginning of
this section, we now consider complete manifolds whose Ricci curvature at
infinity is ``almost nonnegative'' in the following sense. Fix \(p\in M^n\) and
set \(r=\rho(p,x)\). We assume that for some \(r_0>0\) and all \(r\ge r_0\),
\begin{equation}\tag{4.25}
\operatorname{Ric}_M(x)\ge (n-1)\frac{\frac14\pm v^2}{r^2}.
\end{equation}

(For Theorem 4.8 and Theorem 4.9(iii) it suffices to assume that (4.25) holds
only for radial directions from \(p\); however for Theorem 4.9(i) and (ii) this
does not suffice.) By convention we take $v\ge 0$. Observe that the Jacobi
equation
\begin{equation}\tag{4.26}
g''=-\frac{\bigl(\frac14-v^2\bigr)g}{r^2}
\end{equation}
corresponding to sectional curvature \(\bigl(\frac14-v^2\bigr)r^2\) admits the
pair of solutions
\begin{equation}\tag{4.27}
g_{\pm v}(r)=r^{\frac12\pm v},
\end{equation}
and the particular solution
\begin{equation}\tag{4.28}
J_{v,R}=\frac{1}{2v}\Bigl(-r^{\frac12+v}+R^{2v}r^{\frac12-v}\Bigr),
\end{equation}
satisfying \(J_{v,R}(R)=0,\ J'_{v,R}(R)=-1\). For the case
\(\frac14+v^2\), we replace $v$ by $iv$ in (4.27), and obtain the solutions
\begin{equation}\tag{4.29}
r^{\frac12}\cos(v\log r),\qquad r^{\frac12}\sin(v\log r).
\end{equation}
Corresponding to (4.28) we have
\begin{equation}\tag{4.30}
\mathcal G_{iv,R}(r)=\frac{R}{v}\,r^{\frac12}\sin(v\log r/R).
\end{equation}
Note that
\begin{equation}\tag{4.31}
\mathcal G_{iv,R}(\Re^{\pi/v})=\mathcal G_{iv,R}(R)=0.
\end{equation}
This gives the following extension of Myer's theorem.

\begin{theorem}[Asymptotic Ricci lower bound and compactness]
\label{thm:asymptotic-ricci-compactness}
\dcref{Theorem 4.8}
\group{section-4}
\source{cheeger-gromov-taylor:page-0034,cheeger-gromov-taylor:page-0035}
Let \(M^n\) be a complete Riemannian manifold such that for some
$p\in M^n,\, r_0>0,$ and $v>0$ we have
\begin{equation}\tag{4.32}
\operatorname{Ric}_M(x)\ge (n-1)\frac{\frac14+v^2}{r^2}
\end{equation}
for all \(r\ge r_0\). Then \(M^n\) is compact and the diameter $d_p$ from $p$
satisfies
\begin{equation}\tag{4.33}
d_p<e^{\pi/v}r_0.
\end{equation}
\end{theorem}

\begin{proof}[Proof of Theorem 4.8.]
\source{cheeger-gromov-taylor:page-0035}
We use the field $\mathcal{F}_{\gamma|_{[0,e^{\nu}/r_0]}},$ and argue by index
comparison as in the standard proof of Myers' theorem. If
$\gamma|_{[0,e^{\nu}/r_0]}$ is minimal, (4.31) and (4.32) lead to the conclusion
that $\gamma|_{[r_0,e^{\nu}/r_0]}$ contains a pair of conjugate points.
\end{proof}

\begin{remark}[The borderline case in the compactness theorem]
\label{rem:borderline-asymptotic-ricci-compactness}
\dcref{Remark 4.3}
\group{section-4}
\source{cheeger-gromov-taylor:page-0035}
By considering a metric on $\R^n$ which is of the form
\begin{equation}
dr^2 + g(\theta)
\tag{4.34}
\end{equation}
outside some compact set, where $g(\theta)$ is the standard metric on $S^{n-1}$,
we see that $M^n$ need not be compact if $\nu = 0$. However, further refinements
are possible. For example, a similar argument shows that if
\begin{equation}
\Ric_M(x) \ge (n-1)\left\{\frac{1/4}{r^2} + \frac{1/4 + \nu^2}{r^2\log r}\right\},
\tag{4.35}
\end{equation}
then $M^n$ is compact; see also [14] for a sharp result of this type.
\end{remark}

We now consider the case $(\frac14-\nu^2)/r^2$. Observe that for $\nu>0$, the
conditions
\begin{equation}
\nu \le \frac{n+1}{2(n-1)}, \qquad
-\frac{n}{(n-1)^2} \le \frac14-\nu^2, \qquad
0 \le (n-1)\left(\frac12-\nu\right)+1
\tag{4.36}
\end{equation}
are equivalent to one another. Also recall that if $M^n$ is not compact, a
normal geodesic $\gamma:[0,\infty)\to M^n$ is called a ray if each segment of
$\gamma$ is minimal. Parts (i) and (ii) of the following theorem give a lower
bound on the rate of growth of the volume of a ball, which generalizes the
previously mentioned result of Yau and Calabi for the case $\Ric_M\ge 0$.
A novel feature of the argument, as distinct from the case
$\Ric_M\ge (n-1)H$, will be the use of a family of model spaces with metric
\[
dr^2 + J^2_{\nu,R}(r)g,\qquad 0\le r\le R.
\]
Here the metric is expressed in polar coordinates, but the origin is at $R$.
For $\nu<\frac12$, the metric is singular at the boundary $r=0$, while for
$\nu>\frac12$ there is a conjugate point at the boundary. For $\nu=\frac12$
we have the usual flat metric on a ball of radius $R$.

\begin{theorem}[Asymptotic Ricci lower bounds and volume growth]
\label{thm:asymptotic-ricci-volume-growth}
\dcref{Theorem 4.9}
\group{section-4}
\source{cheeger-gromov-taylor:page-0035,cheeger-gromov-taylor:page-0036,cheeger-gromov-taylor:page-0037}
Let $M^n$ be a complete noncompact Riemannian manifold, and let
$\gamma:[0,\infty)\to M^n$ be a ray from some $p\in M^n$.
\begin{enumerate}
\item[(i)] Fix $0\le \nu\le \tfrac12(n+1)/(n-1)$. If $\nu\ge \tfrac12$, assume
that for some $r_0>0$ and all $r>r_0$ we have
\begin{equation}
\tag{4.37}
\Ric_M(x)\ge (n-1)\frac{(\frac14-\nu^2)}{r^2}.
\end{equation}

If $0 \leq \nu < \frac12$, assume that for some $r_1$ such that $r_0 < r_1 <
\frac32 r_0$ and for all $r > r_0$,
\begin{equation}\tag{4.38}
\RicM(x) \geq (n-1)\frac{\left(\frac14-\nu^2\right)}{\bigl(r-2(r_1-r_0)\bigr)^2}.
\end{equation}

If $0 \leq \nu < \frac12 (n+1)/(n-1)$, then there exists $C_{r_1}$,
independent of $M^n$ such that for $r > 2r_1-r_0$
\begin{equation}\tag{4.39}
c_{r_1}V_{(r_1-r_0)}(\gamma(r_1))\,r^{(n-1)(\frac12-\nu)+1} \leq V_{r+(r_1-r_0)}(p).
\end{equation}

\item[(ii)] If $\nu = \frac12 (n+1)/(n-1)$, and (4.37) holds, then for
$r > 2r_1-r_0$
\begin{equation}\tag{4.40}
c_{r_1}V_{(r_1-r_0)}(\gamma(r_1))\log r \leq V_{r+(r_1-r_0)}(p).
\end{equation}

\item[(iii)] For all $\nu \geq 0$, if (4.37) holds (in radial directions), then
\begin{equation}\tag{4.41}
V_r(p) < c_M r^{(n-1)(\frac12+\nu)+1}.
\end{equation}
\end{enumerate}
\end{theorem}

\begin{proof}
\source{cheeger-gromov-taylor:page-0036,cheeger-gromov-taylor:page-0037}
\uses{prop:relative-volume-estimates}
(i) We first consider the case $\nu \geq \frac12$. Fix $x$ with
$\rho(p,x)=r$. Let $\sigma(t)$ be a normal geodesic emanating from $p$, and set
$\rho(\gamma(r),\sigma(t))=s$. We consider $t<r-r_0$ and note that $s>r-t$ by
the triangle inequality. Since $\nu > \frac12$,
\begin{equation}\tag{4.42}
\RicM\sigma(t) \geq (n-1)\frac{\left(\frac14-\nu^2\right)}{s^2} \geq
(n-1)\frac{\left(\frac14-\nu^2\right)}{(r-t)^2}.
\end{equation}

The argument can now be completed as in Proposition 4.1 (iv) except that we
use $J^{n-1}_{\nu,r}$ in place of $\mathscr{Q}^H$. This gives
\begin{equation}\tag{4.43}
\frac{V_{(r_1-r_0)}(\gamma(r_1))}{V_{(r-r_1)}(\gamma(r))}
\leq
\frac{\int_{r_1}^{2r_1-r_0} J^{n-1}_{\nu,r}(t)\,dt}{\int_{2r_1-r_0}^{r} J^{n-1}_{\nu,r}(t)\,dt}
\sim \frac{r^{(n-1)2\nu}}{r^{(n-1)(\frac12+\nu)+1}},
\end{equation}
and (4.39) follows easily.

Now consider the case $0 \leq \nu \leq \frac12$. It will suffice to obtain the
analog of (4.42). Moreover, clearly we can restrict attention to those $\sigma$
which are minimal from $\gamma(r)$ to some $\sigma(t) \in B_{r_1-r_0}(\gamma(r_1))$,
and to $0 \leq t \leq \bar t$. Let $a=\rho(p,\sigma(t))$. Then by the triangle
inequality,
\begin{equation}\tag{4.44}
r_1 + (r_1-r_0) \geq a,
\end{equation}
% SOURCE NOTE: the scanned source is illegible at the marked symbol in the middle terms of (4.45) and (4.46).
\begin{equation}\tag{4.45}
(r-r_1) + (r_1-r_0) \geq t = \text{[illegible]} + t,
\end{equation}
\begin{equation}\tag{4.46}
a + \text{[illegible]} \geq s.
\end{equation}

Adding these and cancelling terms give
\begin{equation}
r + 2(r_1 - r_0) \geq t + s.
\tag{4.47}
\end{equation}

But $r - t \geq r_0$ and so, by the triangle inequality,
\begin{equation}
s \geq r_0.
\tag{4.48}
\end{equation}

Combining (4.47), (4.48) and the assumption $3r_0 > 2r_1$, we obtain
\begin{equation}
r - t \geq s - 2(r_1 - r_0) \geq s - r_0 > 0.
\tag{4.49}
\end{equation}

Thus, since $v < \frac{1}{2}$, we get
\begin{equation}
\Ric_M(x) \geq (n - 1)\frac{\frac{1}{4} - v^2}{(s - 2(r_1 - r_0))^2}
\leq (n - 1)\frac{\frac{1}{4} - v^2}{(r - t)^2}.
\tag{4.50}
\end{equation}

As above, this implies (4.39) for $v < \frac{1}{2}$.

(ii) The argument is as in (i).

(iii) The estimate (4.41) follows from the proof of the Heintze-Karcher
comparison theorem [24] applied to $\partial B_{r_1}\setminus C_1$, where
$r_1 > r_0$ and $C$ is the cut locus of $p$.
\end{proof}

\begin{remark}[Sharpness of the volume-growth estimates]
\label{rem:sharpness-volume-growth-estimates}
\dcref{Remark 4.4}
\group{section-4}
\source{cheeger-gromov-taylor:page-0037}
By considering a metric on $\R^n$ which is of the form
\begin{equation}
dr^2 + r^{1\pm 2v}g
\tag{4.51}
\end{equation}
outside some compact set, we see that the estimates of Theorem 4.9 are sharp.
Also as in Proposition 4.2 we see that the injectivity angle $\tilde{\omega}$
decays at most like $r^{-(n-1)(\frac{1}{2}+v)}$.
\end{remark}

\begin{center}
\textbf{References}
\end{center}

\begin{enumerate}
\item W. Ambrose, \emph{A theorem of Myers}, Duke Math. J. \textbf{24} (1957) 345.
\item M. Berger, \emph{Blaschke's conjecture for spheres}, Appendix D in A. L. Besse,
\emph{Manifolds all of who geodesics are closed}, Ergebnisse Math. und ihrer Grenzgebiete,
Vol. 93, Springer, Berlin, 1978, 236--242.
\item R. Bishop \& B. Crittenden, \emph{Geometry of manifolds}, Academic Press, New York, 1964.
\item E. Bombieri, \emph{Theory of minimal surfaces, and a counter-example to the Bernstein conjecture in
high dimension}, Lecture Notes, Courant Inst., New York, 1970.
\item J. Cheeger, \textit{The relation between the Laplacian and the diameter for manifolds of non-negative curvature}, Arch. Math. 19 (1968) 558--560.
\item \underline{\hspace{2em}}, \textit{Finiteness theorems for riemannian manifolds}, Amer. J. Math. 92 (1970), 61--64.
\item J. Cheeger \& D. G. Ebin, \textit{Comparison theorems in Riemannian geometry}, North Holland, Amsterdam, 1975.
\item J. Cheeger \& M. Taylor, \textit{On the diffraction of waves by conical singularities}, Comm. Pure Appl. Math., to appear.
\item J. Cheeger \& S. T. Yau, \textit{A lower bound for the heat kernel}, Comm. Pure Appl. Math. 34 (1981) 465--480.
\item S. Y. Cheng, P. Li \& S. T. Yau, \textit{On the upper estimate of the heat kernel of a complete riemannian manifold}, Amer. J. Math. 103 (1981) 1021--1063.
\item P. Chernoff, \textit{Essential self-adjointness of powers of generators of hyperbolic equations}, J. Functional Analysis 12 (1973) 401--414.
\item J. Clerc \& E. Stein, \textit{$L^p$ multipliers for noncompact symmetric spaces}, Proc. Nat. Acad. Sci. U.S.A. 71 (1974) 3911--3912.
\item C. Croke, \textit{Some isoperimetric inequalities and eigenvalue estimates}, Ann. Sci. \'{E}cole Norm. Sup. 13 (1980) 419--435.
\item B. V. Dekster \& I. Kupka, \textit{Asymptotics of curvature in a space of positive curvature}, J. Differential Geometry 15 (1980) 553--568.
\item G. de Rham, \textit{Varieties differentiables}, Publ. Inst. Math. Univ. Nancago III, Hermann, Paris, 1960.
\item H. Donnelly, \textit{Spectral geometry for certain noncompact riemannian manifolds}, Math. Z. 169 (1979) 63--76.
\item H. Federer \& W. Flemming, \textit{Normal integral currents} Ann. of Math. 72 (1960) 458--520.
\item G. Galloway, \textit{Compactness criteria for Riemannian manifolds}, Proc. Amer. Math. Soc., 84 (1982) 106--110.
\item L. G\aa rding, \textit{Vecteurs analytiques dans les representations des groupes de Lie}, Bull. Soc. Math. France 188 (1960) 73--93.
\item D. Gilbarg \& M. Trudinger, \textit{Elliptic partial differential equations of second order}, Springer, Berlin, 1977.
\item M. Gromov, \textit{Almost flat manifolds}, J. Differential Geometry 13 (1978) 231--241.
\item \underline{\hspace{2em}}, \textit{Curvature, diameter, and Betti numbers}, Comment. Math. Helv. 56 (1981) 179--197.
\item \underline{\hspace{2em}}, \textit{Paul Levy's isoperimetric inequality}, Inst. Hautes \'{E}tudes Sci. Publ. Math., to appear.
\item E. Heintze \& H. Karcher, \textit{A general comparison theorem with applications to volume estimates for submanifolds}, Ann. Sci. \'{E}cole Norm. Sup. 11 (1978) 451--470.
\item J. Kazdan, \textit{An inequality arising in geometry}, Appendix E, in A. L. Besse, \textit{Manifolds all of whose geodesics are closed}, Ergebnisse Math. und ihrer Grenzgebiete, Vol. 93, Springer, Berlin, 1978, 243--246.
\item N. N. Lebedev, \textit{Special functions and their applications}, Dover, New York, 1972.
\item J. Lions \& E. Magenes, \textit{Non-homogeneous boundary value problems and applications}, Vol. 3, Springer, New York, 1972.
\item V. G. Maz'ya, \textit{Classes of domains and imbedding theorems for function spaces}, Dokl. Akad. Nauk SSSR 133 (1960) 527--530; English Translation, Soviet Math. Dokl. 1 (1960) 882--885.
\item J. Milnor, \textit{A note on curvature and fundamental group}, J. Differential Geometry 2 (1968) 1--7.
\item J. Moser, \textit{A Harnack inequality for parabolic differential equations}, Comm. Pure Appl. Math. 17 (1964) 101--134.
\item R. Stanton \& P. Tomas, \textit{Expansions for spherical functions on noncompact symmetric spaces}, Acta. Math. 140 (1978) 251--276.
\item M. Taylor, \textit{Pseudodifferential operators}, Princeton Univ. Press, Princeton, 1981.
\end{enumerate}

\begin{center}
\textsc{State University of New York, Stony Brook}
\end{center}
